\documentclass[UTF8,zihao=-4,openany]{ctexrep}

\usepackage[a4paper,top=2.6cm,bottom=2.6cm,left=2.8cm,right=2.6cm]{geometry}
\usepackage{amsmath,amssymb,bm}
\usepackage{booktabs,longtable,tabularx,array,multirow}
\usepackage{graphicx,float,placeins}
\usepackage{xcolor}
\usepackage{enumitem}
\usepackage{fancyhdr}
\usepackage{hyperref}
\usepackage{caption}
\usepackage{subcaption}
\usepackage{titlesec}
\usepackage{setspace}

% 参赛信息：提交前请替换“待填写”内容。
\newcommand{\ApplicantUnit}{待填写}
\newcommand{\ApplicantName}{待填写}
\newcommand{\TeamMembers}{待填写}
\newcommand{\AdvisorNames}{待填写}
\newcommand{\ContactPhone}{待填写}
\newcommand{\ReportDate}{2026年9月13日}

% 已验证数据规模。
\newcommand{\PlateCount}{109}
\newcommand{\PartCount}{910}
\newcommand{\KitGroupCount}{21}

% 当前代码版本复现实验结果；由复现实验完成后更新。
\newcommand{\BaseCmax}{63.44}
\newcommand{\OptCmax}{59.22}
\newcommand{\BalancedCmax}{60.55}
\newcommand{\BaseKitSpan}{13.89}
\newcommand{\OptKitSpan}{11.40}
\newcommand{\BalancedKitSpan}{10.99}
\newcommand{\BaseLoadDiff}{0.63}
\newcommand{\OptLoadDiff}{2.61}
\newcommand{\BalancedLoadDiff}{0.01}
\newcommand{\BaseThroughput}{13.75}
\newcommand{\OptThroughput}{14.73}
\newcommand{\BalancedThroughput}{14.40}
\newcommand{\BaseCutUtil}{80.33}
\newcommand{\OptCutUtil}{86.05}
\newcommand{\BalancedCutUtil}{84.15}
\newcommand{\BaseBevelWorkstationPeak}{114.29}
\newcommand{\OptBevelWorkstationPeak}{57.14}
\newcommand{\BalancedBevelWorkstationPeak}{100.00}
\newcommand{\DeadlockCount}{0}


\definecolor{steelblue}{HTML}{1F4E79}
\definecolor{lightsteel}{HTML}{D9EAF7}
\definecolor{linegray}{HTML}{D9D9D9}
\definecolor{textgray}{HTML}{555555}

\hypersetup{
  colorlinks=true,
  linkcolor=steelblue,
  citecolor=steelblue,
  urlcolor=steelblue,
  pdfauthor={\ApplicantUnit},
  pdftitle={船体加工车间智能排产与齐套配盘优化调度技术方案报告}
}

\setstretch{1.32}
\setlength{\headheight}{13pt}
\setlength{\parindent}{2em}
\setlength{\parskip}{0.15em}
\setlist{nosep,leftmargin=2.2em}
\captionsetup{font=small,labelsep=quad}
\renewcommand{\arraystretch}{1.22}
\setlength{\tabcolsep}{5pt}

\titleformat{\chapter}{\centering\bfseries\zihao{2}\color{black}}{第\chinese{chapter}章}{1em}{}
\titleformat{\section}{\bfseries\zihao{3}\color{black}}{\thesection}{0.8em}{}
\titleformat{\subsection}{\bfseries\zihao{4}\color{black}}{\thesubsection}{0.8em}{}

\pagestyle{fancy}
\fancyhf{}
\fancyhead[L]{\small 船体加工车间智能排产与齐套配盘优化调度}
\fancyhead[R]{\small 技术方案报告}
\fancyfoot[C]{\thepage}
\renewcommand{\headrulewidth}{0.4pt}

\newcolumntype{Y}{>{\raggedright\arraybackslash}X}
\newcommand{\code}[1]{\texttt{#1}}
\newcommand{\important}[1]{\textbf{#1}}

% 最终版举证图片；缺图时由 LaTeX 直接报错，避免占位内容进入提交稿。
\newcommand{\evidencefigure}[3][0.94\textwidth]{%
  \begin{figure}[H]
    \centering
    \includegraphics[width=#1,height=0.68\textheight,keepaspectratio]{#2}%
    \caption{#3}
  \end{figure}
}

\begin{document}

\begin{titlepage}
  \centering
  \vspace*{2.0cm}
  {\zihao{1}\bfseries 船体加工车间智能排产与齐套配盘优化调度\par}
  \vspace{0.8cm}
  {\zihao{2}\bfseries 技术方案报告\par}
  \vspace{0.6cm}
  {\zihao{4}\color{steelblue}题目编号 CS-202625\par}
  \vfill
  \renewcommand{\arraystretch}{1.5}
  \begin{tabular}{>{\bfseries}r l}
    申报单位： & \ApplicantUnit \\
    申报人： & \ApplicantName \\
    团队成员： & \TeamMembers \\
    指导教师： & \AdvisorNames \\
    联系电话： & \ContactPhone \\
    报告日期： & \ReportDate \\
  \end{tabular}
  \vfill
  {\small 本报告对应同目录提交的算法源代码、仿真测试程序及实验举证材料。\par}
\end{titlepage}

\pagenumbering{Roman}
\chapter*{摘要}
\addcontentsline{toc}{chapter}{摘要}

本方案针对船体加工车间多品种、小批量、多工序耦合和齐套优先的排产问题，建立覆盖钢板上料、N2/N5双工位切割、桁架分拣、小件自动打磨与坡口、大件人工打磨与坡口、AGV及天车转运、有限缓存和分段齐套的离散事件仿真模型。系统以钢板加工顺序和切割机分配为主要决策，把设备资格、资源互斥、工序前后序、料框分类、缓存容量和齐套门控统一到同一事件时间轴中。

求解层以可稳定复现的SA+Tabu为当前主验证路径，同时保留GA+LNS和DQN种子增强NSGA-II接口。各算法调用同一仿真评价器，以总完工时间、加权平均齐套跨度、切割负载差、等待时间和设备利用率为核心指标。当前数据包含\PlateCount 张钢板、\PartCount 个零件和\KitGroupCount 个分段优先级齐套组。最新代码版本的正式复现实验数值由\code{report\_data.tex}统一维护，避免正文、图表和程序输出出现多套口径。

系统实现FastAPI服务端与React可视化界面，支持附件数据校验、参数配置、排产计算、甘特图、齐套分析、设备利用率、缓存监控、运行历史和设备故障后的剩余任务重排。对于题目未提供的缓存停留、物流和现场容量参数，本方案集中置于\code{ModelConfig}中，并将其标明为待现场校准参数，不作为企业实测结论。

固定随机种子20260723、260次局部搜索预算下，产能优先方案将总完工时间由FIFO的63.44小时降至59.22小时，加权平均齐套跨度由13.89小时降至11.40小时，整体产能提高至14.73张/班；三指标平衡方案将齐套跨度进一步降至10.99小时，并把N2/N5切割负载差降至0.01小时。两条轨道均未出现缓存溢出或死锁警告。

\noindent\textbf{关键词：}船体加工；智能排产；齐套配盘；离散事件仿真；SA+Tabu；NSGA-II；动态重调度

\tableofcontents
\clearpage
\pagenumbering{arabic}

\chapter{项目背景与赛题响应}

\section{问题背景}

船体加工车间的交付对象不是单张钢板，而是由多个钢板、不同类型零件和多道工序共同构成的分段齐套集合。若同一分段中少数关键零件延迟，已经完工的其他零件仍需占用料框和缓存区等待，形成在制品积压。与此同时，N2与N5切割机具有不同设备资格和负载状态，下游打磨、坡口、AGV及天车又会反向制约切割顺序。因此，单纯缩短切割时间或采用先到先服务规则，无法同时获得高产能和短齐套周期。该问题属于带异构资源、前后序与有限缓存约束的组合调度问题\cite{pinedo}。

\section{方案贡献}

本方案的核心贡献包括：以分段齐套而非单板完工为评价中心；把切割、下游加工、转运与有限缓存置于同一离散事件时间轴；以同一仿真评价器支撑SA+Tabu、GA+LNS和NSGA-II，避免算法之间评价口径漂移；并在设备故障发生后冻结已执行任务，仅对剩余任务重排，使静态优化结果能够延伸到现场动态响应。

\section{赛题要求与本方案对应}

\begin{table}[H]
\centering
\caption{赛题要求与系统实现对应关系}
\begin{tabularx}{\textwidth}{>{\bfseries}p{0.23\textwidth}YY}
\toprule
赛题要求 & 本方案实现 & 可验证证据 \\
\midrule
多工序和异构设备建模 & 统一建模切割、分拣、打磨、坡口、转运、缓存和齐套事件 & 数学模型、工序事件表、资源利用率 \\
有限缓存与齐套配盘 & 按分段、优先级、坡口属性、零件类型和来源机台组框，并记录容量占用 & 齐套结果、缓存时序、峰值占用率 \\
设备负载均衡 & 钢板顺序与N2/N5分配联合评价，负载差进入目标函数 & 双机甘特图、N2/N5利用率 \\
多目标动态优化 & 同时评价总完工时间、齐套跨度、负载差和等待时间 & FIFO与优化方案对比、Pareto接口 \\
设备故障和订单扰动 & 冻结已完成/不受影响任务，对剩余任务在故障日历上重排 & 动态重排接口和前端重调度页 \\
完整仿真系统 & FastAPI后端、React前端、CSV/JSON/PNG导出 & 可运行程序、运行说明和演示截图 \\
\bottomrule
\end{tabularx}
\end{table}

\section{参考指标与评价口径}

题目给出的参考指标为整体产能13张板/班、总完工时间67.25小时、切割机稼动率61\%和平均齐套时长12.7小时。程序将这些值作为对照口径，同时使用由当前数据自动计算的理论下界进行归一化，避免数据规模变化后仍机械套用固定数值。报告只引用同一代码版本、同一输入和同一随机种子生成的结果。

\chapter{系统总体方案}

\section{程序的实际含义}

该程序本质上是一套\important{面向船体钢板加工的数字化排产实验平台}。它不是生产线控制系统，也不直接驱动PLC；其作用是读取钢板与零件数据，构造候选加工顺序，在离散事件环境中推演每个资源的占用和释放，再由优化算法反复搜索更好的排产方案，最后以甘特图、齐套图和KPI表的形式供计划人员比较。离散事件建模与实验验证遵循仿真建模的一般方法\cite{law}。

\begin{figure}[H]
\centering
\fbox{\begin{minipage}{0.92\textwidth}
\centering
\textbf{输入数据层}\\
附件2钢板与零件数据；附件3工艺速度和用时规则\\[0.6em]
$\Downarrow$\\[0.6em]
\textbf{模型与仿真层}\\
数据校验 $\rightarrow$ 工艺时间计算 $\rightarrow$ 资源池与事件队列 $\rightarrow$ 齐套及缓存状态\\[0.6em]
$\Downarrow$\\[0.6em]
\textbf{优化求解层}\\
SA+Tabu；GA+LNS；NSGA-II（可选DQN机器分配种子）\\[0.6em]
$\Downarrow$\\[0.6em]
\textbf{服务与展示层}\\
FastAPI接口 $\rightarrow$ React页面 $\rightarrow$ CSV、JSON、PNG和Word导出
\end{minipage}}
\caption{系统数据流和功能分层}
\end{figure}

\section{软件结构}

\small
\begin{longtable}{p{0.36\textwidth}p{0.56\textwidth}}
\caption{核心代码模块}\label{tab:code-map}\\
\toprule
文件或目录 & 职责 \\
\midrule
\endfirsthead
\toprule
文件或目录 & 职责 \\
\midrule
\endhead
\code{steel\_schedule\_model.py} & 数据校验、工艺参数、资源池、离散事件仿真、KPI和图表输出 \\
\code{improved\_optimizer.py} & SA+Tabu、多初始策略、邻域搜索、缓存评价和并行入口 \\
\code{ga\_lns\_optimizer.py} & 遗传编码、交叉变异、破坏修复和档案维护 \\
\code{dqn\_machine\_model.py} & DQN机器分配网络、特征域检查、权重加载和启发式回退 \\
\code{dqn\_nsga2\_optimizer.py} & 钢板顺序与机器码联合编码、非支配排序和拥挤度筛选 \\
\code{drl\_optimizer.py} & 轻量Q-learning算子选择实验模块 \\
\code{backend/main.py} & 数据上传、运行管理、结果读取、报告导出、重调度和鲁棒性接口 \\
\code{frontend/src} & 参数配置、结果概览、甘特图、齐套分析、利用率和缓存监控 \\
\bottomrule
\end{longtable}
\normalsize

\section{输入输出}

附件2包含钢板表和零件表。程序检查主键重复、钢板与零件关联、零件数量、小件/大件数量和V坡长度汇总关系，并保留打磨长度缺失清单。附件3提供厚度相关的直切、V\&Y坡和X\&K坡速度，非精确厚度使用线性插值。

每次完整运行输出钢板排程、零件完成时刻、齐套组、工序事件、方案对比、缓存时序和汇总指标。Web端将这些数据转换为计划员可直接查看的图表和历史记录。

\chapter{生产过程与建模假设}

\section{资源边界}

当前正式场景固定为1台人工天车、N2/N5两台双工位切割机、2套双臂分拣桁架、2套小件自由边打磨设备、1台自动小件坡口设备、1个人工坡口工位和2台AGV。N5仅接收板厚不超过36\,mm且板宽不超过4500\,mm的钢板；不满足资格条件的钢板必须分配至N2。

\section{工艺路线}

小件在切割后经过桁架分拣、自动打磨、按坡口属性分流、码垛和AGV转运。坡口件还需进入自动坡口工位；当同一分段的相关优先级组满足到齐条件后，料框由AGV转入齐套区。大件经过胎架自由边打磨，需坡口时由天车运至人工坡口工位，随后进入成品放置区。

\section{料框与齐套定义}

小件料框键定义为
\begin{equation}
 b(i)=\bigl(s_i,q_i,e_i,t_i,m_i\bigr),
\end{equation}
其中$s_i$为分段号，$q_i$为齐套优先级，$e_i$表示是否需要坡口，$t_i$为零件类型，$m_i$为来源切割机。坡口件与非坡口件不混装。程序首先判断各优先级组是否到齐，再按分段执行最终齐套转运。

组$g$的齐套跨度定义为
\begin{equation}
 K_g=\max_{i\in I_g}a_i-\min_{i\in I_g}a_i,
\end{equation}
其中$a_i$为零件到达中间缓存或坡口工作站的时间。加权平均齐套跨度为
\begin{equation}
 \bar K=\frac{\sum_{g\in G}|I_g|K_g}{\sum_{g\in G}|I_g|}.
\end{equation}

\section{参数边界}

凡题目未给出的AGV单次运输时间、齐套区停留时间、部分缓存容量和空车返回系数，均集中在\code{ModelConfig}中。正式答辩时必须把这些参数称为“模型参数”或“待现场校准参数”，不能称为企业实测值。

\chapter{数学模型}

\section{集合和参数}

设$P$为钢板集合，$I$为零件集合，$M=\{\mathrm{N2},\mathrm{N5}\}$为切割机集合，$R$为下游资源集合，$G$为齐套组集合，$H$为缓存集合。$I_p$表示钢板$p$上的零件，$I_g$表示齐套组$g$包含的零件。$d_p^{\mathrm{cut}}$为钢板切割工时，$d_{ir}$为零件$i$在资源$r$上的加工时间，$B_h$为缓存$h$的容量。

\section{切割时间}

程序采用附件3的厚度相关速度表。分项切割时间为
\begin{equation}
 d_p^{\mathrm{cut}}=
 \frac{L_p^{\mathrm{straight}}/v^{\mathrm{straight}}(\theta_p)
 +L_p^{V}/v^{VY}(\theta_p)}{0.8\times0.6}
 +\frac{L_p^{\mathrm{empty}}}{v^{\mathrm{rapid}}}
 +\frac{L_p^{\mathrm{mark}}}{v^{\mathrm{mark}}}
 +d_p^{\mathrm{remainder}}.
\end{equation}
穿孔影响已折算进题目给定系数，不再重复累计。双工位切割场景中，残材处理与另一工位加工重叠，默认不作为切割头串行占用时间。

\section{决策变量}

用$x_{pm}\in\{0,1\}$表示钢板$p$是否分配给切割机$m$，用排列$\pi=(p_1,p_2,\ldots,p_{|P|})$表示钢板加工顺序。$S_{ir}$和$C_{ir}$分别表示任务$i$在资源$r$上的开始和完成时刻。基本分配约束为
\begin{equation}
 \sum_{m\in M}x_{pm}=1,\qquad \forall p\in P,
\end{equation}
并对N5资格不满足的钢板强制$x_{p,\mathrm{N5}}=0$。

\section{工序和资源约束}

对零件$i$的相邻工序$k$与$k+1$，有
\begin{equation}
 S_{i,k+1}\ge C_{ik}+\tau_{i,k\rightarrow k+1},
\end{equation}
其中$\tau$为必要转运时间。同一资源上的任意两个任务$a,b$必须满足析取约束
\begin{equation}
 C_{ar}\le S_{br}\quad\text{或}\quad C_{br}\le S_{ar}.
\end{equation}
资源故障区间记为$\Omega_r$，任务占用区间不得与$\Omega_r$相交。缓存占用写为
\begin{equation}
 0\le Q_h(t)\le B_h,\qquad h\in H.\label{eq:buffer-capacity}
\end{equation}
机旁码垛区和齐套缓存采用事件延迟门控；半框缓存、坡口缓存和坡口工作站通过峰值占用率参与可行性评价。任一容量超限或死锁候选解都会被施加高额惩罚，并在最终选解阶段优先排除，从而保证正式输出方案满足式~\eqref{eq:buffer-capacity}的容量边界。

\section{优化目标}

总完工时间为
\begin{equation}
 C_{\max}=\max\left\{\max_{g\in G}C_g^{\mathrm{kit}},\ \max_{i\in I_{\mathrm{large}}}C_i\right\}.
\end{equation}
切割负载差为
\begin{equation}
 D_{\mathrm{load}}=\max_{m\in M}\sum_{p\in P}x_{pm}d_p^{\mathrm{cut}}
 -\min_{m\in M}\sum_{p\in P}x_{pm}d_p^{\mathrm{cut}}.
\end{equation}
当前主路径采用产能优先的有界辅助目标
\begin{equation}
 J=\frac{C_{\max}}{LB}+\varepsilon
 \left(w_K n_K+w_L n_L+w_W n_W\right),
\end{equation}
其中$LB$为理论下界，$n_K,n_L,n_W\in[0,1]$分别为齐套跨度、负载差和等待时间的归一化量。$\varepsilon$限制辅助指标对产能目标的扰动范围。

NSGA-II接口直接使用四维目标向量
\begin{equation}
 \bm f=\left(\frac{C_{\max}}{LB},\frac{\bar K}{\bar K_{\mathrm{FIFO}}},
 \frac{D_{\mathrm{load}}}{D_{\mathrm{FIFO}}},\frac{W}{W_{\mathrm{FIFO}}}\right).
\end{equation}

\chapter{优化算法设计}

\section{SA与Tabu混合搜索}

当前可完整复现的主算法以FIFO、分段优先、短工时优先、长工时优先和天车感知等规则构造多个初始解。搜索过程中交替使用交换、插入、分块逆序、分段聚类、齐套聚类和跨机负载修复等邻域。模拟退火以概率
\begin{equation}
 P(\text{accept})=\min\left(1,\exp\left[-\frac{J'-J}{T}\right]\right)
\end{equation}
接受暂时变差的解\cite{sa}，Tabu表记录近期序列摘要，降低循环搜索概率\cite{tabu}。所有候选解都调用同一离散事件仿真器计算指标。

\section{GA与LNS}

GA+LNS以钢板排列为主要染色体，采用顺序交叉和排列变异保持钢板不重复。LNS按随机、分段或瓶颈特征破坏一部分序列，再根据齐套优先级和当前资源负载修复。该算法适合需要扩大搜索范围的离线实验，但正式结果必须使用与主实验相同的输入、参数和评价口径。

\section{DQN种子与NSGA-II}

DQN模块按16维状态输入，经$16\rightarrow128\rightarrow64\rightarrow3$网络输出机器分配动作，在N2/N5正式场景中仅使用两台设备对应动作。它只为NSGA-II提供可选机器码种子，不替代完整排产仿真；当PyTorch、权重文件或特征范围校验不通过时，系统自动回退到启发式种子。本文正式实验采用SA+Tabu，不将DQN结果作为性能结论依据。

NSGA-II联合演化钢板顺序$OS$和机器码$MS$，使用非支配排序和拥挤距离保留多目标前沿。一次搜索可从前沿中分别选择产能优先解和综合平衡解。

DQN网络结构参考深度Q网络的基本形式\cite{dqn}，非支配排序和拥挤距离采用经典NSGA-II定义\cite{nsga2}。

\section{算法选择原则}

\begin{table}[H]
\centering
\caption{算法适用边界}
\begin{tabularx}{\textwidth}{p{0.20\textwidth}YY}
\toprule
算法 & 适用场景 & 当前交付状态 \\
\midrule
SA+Tabu & 中等规模、需要快速稳定得到单一方案 & 当前正式复现主路径 \\
GA+LNS & 需要更强全局多样性的离线搜索 & 已实现，需单独固定参数复现 \\
NSGA-II & 同时比较产能、齐套、负载和等待 & 已实现多目标前沿 \\
DQN机器分配 & 有可信训练数据和随包权重时改善初始机器码 & 可选扩展；当前正式名称不宣称DQN已启用 \\
DRL算子选择 & 用历史搜索轨迹学习邻域与权重选择 & 实验模块，不作为当前主结果依据 \\
\bottomrule
\end{tabularx}
\end{table}

\chapter{离散事件仿真与动态响应}

\section{状态模型}

仿真状态包含切割头和双工位释放时间、下游资源池、天车时序、AGV位置与区域释放时间、料框组成、缓存进入/离开事件、齐套组到达状态和故障日历。事件以时间和稳定序号排序，保证同一输入和同一随机种子下结果可追踪。

\section{事件推进}

仿真按以下顺序推进：数据校验；钢板上料和双工位切割；小件与大件分流；分拣、打磨及坡口加工；料框聚合和转运；齐套条件判断；最终齐套转运；KPI与文件输出。与只统计总工时的静态模型相比，该过程能够显式发现资源等待、物流冲突和缓存压力。

\section{设备故障重排}

当前接口支持\code{machine\_failure}、\code{rush\_order}和\code{reoptimize}三类请求。设备故障场景在决策时刻冻结不受影响或已经完成的钢板，把故障机台在故障窗口内的任务延后，并对剩余钢板进行短时SA+Tabu搜索。重排完成后重新运行全链路仿真，返回新的甘特图、工序事件、设备利用率和缓存时序。

\code{rush\_order}接口用于将当前数据集中已有钢板在剩余任务队列中优先前移，功能口径为“紧急任务优先调整”。任意新增订单的主数据合并、零件关联校验和版本管理不属于本次交付范围。

\chapter{仿真实验设计与结果}

\section{实验环境}

\begin{table}[H]
\centering
\caption{复现实验环境}
\begin{tabularx}{\textwidth}{p{0.25\textwidth}Y}
\toprule
项目 & 配置 \\
\midrule
操作系统 & Windows \\
Python & 3.11，64位 \\
前端 & React 19、TypeScript 6、Vite 8、Ant Design 6、ECharts 6 \\
后端 & FastAPI、pandas、NumPy、openpyxl、Pillow、matplotlib \\
输入 & 附件2钢板零件数据、附件3工艺用时计算表 \\
随机种子 & 20260723 \\
主验证算法 & SA+Tabu，线性产能优先评价，260次局部搜索预算 \\
\bottomrule
\end{tabularx}
\end{table}

\section{数据校验}

官方数据共包含\PlateCount 张钢板和\PartCount 个零件，形成\KitGroupCount 个“分段+优先级”齐套组。钢板与零件关联、零件数量、小件/大件数量和V坡汇总检查均通过。原始数据中有2个零件缺少打磨长度，程序将其记为数据告警并按0参与本次仿真；该处理仅用于保证赛题数据可计算，不代表现场确认值。

\section{主实验结果}

\begin{table}[H]
\centering
\caption{FIFO基线与当前版本两条优化轨道结果}
\small
\begin{tabularx}{\textwidth}{Ycccc}
\toprule
指标 & FIFO基线 & 产能优先 & 三指标平衡 & 赛题参考值 \\
\midrule
总完工时间（h） & \BaseCmax & \OptCmax & \BalancedCmax & 67.25 \\
加权平均齐套跨度（h） & \BaseKitSpan & \OptKitSpan & \BalancedKitSpan & 12.7 \\
切割负载差（h） & \BaseLoadDiff & \OptLoadDiff & \BalancedLoadDiff & -- \\
整体产能（张板/班） & \BaseThroughput & \OptThroughput & \BalancedThroughput & 13 \\
切割机综合稼动率 & \BaseCutUtil\% & \OptCutUtil\% & \BalancedCutUtil\% & 61\% \\
坡口工作站峰值占用率 & \BaseBevelWorkstationPeak\% & \OptBevelWorkstationPeak\% & \BalancedBevelWorkstationPeak\% & $\le100\%$ \\
死锁警告数 & -- & \DeadlockCount & \DeadlockCount & 0 \\
\bottomrule
\end{tabularx}
\end{table}

本表的数值只从当前代码版本复现实验读取。相较FIFO，产能优先轨道将总完工时间缩短约6.65\%，加权平均齐套跨度缩短约17.94\%，整体产能提高约7.09\%；三指标平衡轨道将总完工时间缩短约4.55\%，齐套跨度缩短约20.91\%，并将N2/N5负载差降低约98.37\%。产能优先轨道取得最低总完工时间和最高产能，但切割负载差较大；三指标平衡轨道几乎消除N2/N5负载差，并把坡口工作站峰值控制在容量边界。两条优化轨道的缓存峰值均不超过100\%，死锁警告均为0。若后续更换参数或修改约束，必须重新生成全部指标与图片，不能将不同运行的最优值拼接成一组结果。

\evidencefigure{assets/cutting_gantt.png}{当前版本优化排程的N2/N5切割甘特图}
\evidencefigure{assets/kit_span.png}{各分段优先级组的齐套跨度}
\evidencefigure{assets/resource_utilisation.png}{下游资源利用率}

\section{Web系统举证材料}

以下静态页面截图取自同一组正式运行，保留了当前输入、算法轨道、运行记录和主要指标；动态重调度截图以该排程为基础注入N5故障。它们依次证明数据与参数可配置、排程结果可解释、资源与齐套状态可追踪，以及异常发生后能够在线重排。

\evidencefigure{assets/ui_data_upload.png}{附件2与附件3数据上传界面}
\evidencefigure{assets/ui_parameter_config.png}{优化算法、评价函数与目标权重配置}
\evidencefigure{assets/ui_result_overview.png}{FIFO基线与齐套感知优化方案关键指标对比}
\evidencefigure{assets/ui_gantt.png}{全资源排产甘特图与工序时间轴}
\evidencefigure{assets/ui_kit_analysis.png}{各齐套组首件到达至齐套完成区间}
\evidencefigure{assets/ui_kit_dashboard.png}{齐套配盘看板、瓶颈分段与分段进度}
\evidencefigure{assets/ui_resource_utilization.png}{设备资源利用率统计}
\evidencefigure{assets/ui_buffer_monitor.png}{四级缓存时序、容量阈值与溢出状态监控}
\evidencefigure{assets/ui_reschedule_result.png}{N5中故障场景的动态重排结果}

\section{动态实验设计}

正式举证采用N5在第12小时发生、持续3小时的中故障场景。系统冻结18张已执行或不受影响的钢板，对其余91张钢板重新优化；本次重排响应时间为5.49秒，新总完工时间为64.8小时，加权齐套跨度为13.36小时，整体产能为13.5张/班，平均稼动率为78.6\%。截图中的故障参数、响应时间和重排后KPI来自同一次运行，可据此复核动态重调度链路。

\chapter{工程实现与复现方法}

\section{启动方法}

在Windows环境下可双击\code{start.bat}启动。脚本完成依赖检查、前端构建和FastAPI服务启动后，浏览器访问\url{http://127.0.0.1:8000}。命令行基础模型可按以下方式运行：

\begin{verbatim}
python steel_schedule_model.py \
  --input "附件2：钢板零件数据.xlsx" \
  --speed-table "附件3：工艺用时计算表.xlsx" \
  --output outputs
\end{verbatim}

\section{操作流程}

\begin{enumerate}
  \item 上传附件2，并按需上传附件3；确认数据校验全部通过。
  \item 保持N2/N5、资源数量和缓存容量与正式实验配置一致。
  \item 选择算法、评价函数、迭代预算和随机种子后开始计算。
  \item 查看结果概览、甘特图、齐套分析、设备利用率和缓存监控。
  \item 导出CSV、PNG和运行报告，并记录运行ID及参数。
  \item 在重调度页配置故障机台、故障时刻和持续时间，执行动态验证。
\end{enumerate}

\section{可复现性控制}

本提交包固定附件2、附件3、随机种子20260723、260次局部搜索预算和当前并行配置，并保存两条优化轨道的\code{summary.json}、CSV与PNG。所有正文指标均由这两份汇总文件生成。复现时应保持上述配置一致；如修改模型约束或参数，应把新结果保存到独立目录并同步更新报告数据，避免跨版本拼接指标。

\chapter{工程边界与交付结论}

\section{能力边界与验证状态}

\begin{longtable}{p{0.22\textwidth}p{0.34\textwidth}p{0.34\textwidth}}
\caption{系统能力边界与验证证据}\label{tab:delivery-boundary}\\
\toprule
能力项 & 本次交付口径 & 验证证据 \\
\midrule
\endfirsthead
\toprule
能力项 & 本次交付口径 & 验证证据 \\
\midrule
\endhead
静态智能排产 & 联合决定109张钢板的加工顺序与N2/N5分配，并推进完整下游工序 & 固定种子汇总文件、排程CSV、甘特图与FIFO对比 \\
齐套配盘优化 & 以“分段+优先级”组为齐套单元，评价加权平均与最大齐套跨度 & 齐套组CSV、齐套跨度图和齐套看板 \\
有限缓存约束 & 机旁码垛与齐套缓存采用事件门控，其余缓存进入可行性评价 & 缓存时序、峰值占用率和0项死锁告警 \\
动态重调度 & 支持机器故障、已有任务优先调整和全局重优化 & N5中故障场景、冻结/重排数量及5.49秒响应记录 \\
算法交付 & SA+Tabu为正式复现主路径；GA+LNS与NSGA-II作为可选求解器 & 统一仿真评价接口与源代码；正式结论不依赖DQN权重 \\
系统定位 & 面向计划员的排产决策与仿真验证平台，不直接控制PLC或替代MES & FastAPI接口、React界面、导出文件和运行说明 \\
数据边界 & 两个缺失打磨长度按0处理并显式告警；未知现场参数集中配置 & 数据校验摘要与\code{ModelConfig}参数表 \\
\bottomrule
\end{longtable}

\section{技术创新与应用价值}

\begin{enumerate}
  \item \textbf{齐套目标前置。} 将分段齐套跨度直接纳入排产评价，使优化对象从“单张钢板尽快完成”扩展为“下游装配所需零件尽早成套”。
  \item \textbf{生产物流一体化仿真。} 在同一事件时间轴内表达切割、分拣、打磨、坡口、AGV、天车和缓存，能够解释局部加速为何可能转化为下游拥堵。
  \item \textbf{可行性优先的多轨道优化。} 所有算法共享同一仿真评价器，容量超限与死锁解被统一惩罚；产能优先轨道和三指标平衡轨道可服务不同计划目标。
  \item \textbf{面向扰动的滚动重排。} 故障发生后保留已执行决策，仅优化受影响的剩余任务，在控制计划波动的同时给出秒级响应和新的KPI结果。
\end{enumerate}

\section{结论}

本方案已形成从附件数据校验、工艺时间计算、钢板排序与机器分配，到全链路离散事件仿真、齐套判断、有限缓存监控、KPI导出和设备故障重排的完整闭环。产能优先方案相较FIFO将总完工时间缩短6.65\%、加权平均齐套跨度缩短17.94\%、整体产能提高7.09\%；三指标平衡方案在保持14.40张/班产能的同时，将齐套跨度缩短20.91\%，并将切割负载差降至0.01小时。最终交付物包括技术报告、LaTeX源文件、算法与仿真程序、官方输入以及可复核的JSON、CSV、PNG和Web截图，能够支持现场演示、结果追溯与后续参数校准。

\appendix
\chapter{主要参数表}

\begin{longtable}{p{0.32\textwidth}p{0.20\textwidth}p{0.38\textwidth}}
\toprule
参数 & 当前值 & 来源或说明 \\
\midrule
直切速度 & 按厚度查表 & 附件3；无表时回退1700\,mm/min \\
V\&Y坡速度 & 按厚度查表 & 附件3；线性插值 \\
空行程/划线速度 & 24000\,mm/min & 题目工艺说明 \\
小件打磨速度 & 2520\,mm/min & 42\,mm/s \\
大件打磨速度 & 1950\,mm/min & 1.95\,m/min \\
人工坡口速度 & 250\,mm/min & 题目工艺说明 \\
N2/N5码垛容量 & 10/18框 & 题目给定 \\
坡口缓存 & 3框 & 题目给定 \\
半框缓存 & 14框 & 模型参数，待现场确认 \\
齐套缓存 & 16框 & 模型参数，待现场确认 \\
齐套停留时间 & 60\,min & 模型参数，待现场确认 \\
\bottomrule
\end{longtable}

\chapter{提交材料对应关系}

\begin{table}[H]
\centering
\begin{tabularx}{\textwidth}{p{0.28\textwidth}Y}
\toprule
提交物 & 提交包位置 \\
\midrule
技术报告与LaTeX源文件 & \code{01\_技术报告} \\
算法源代码与仿真程序 & \code{02\_算法与仿真程序} \\
赛题文件及附件2/3 & \code{03\_赛题与输入数据} \\
固定实验结果 & \code{04\_仿真实验举证材料/01--02} \\
Web系统截图 & \code{04\_仿真实验举证材料/03\_Web系统截图} \\
运行与提交说明 & 提交包根目录\code{README\_提交说明.md} \\
\bottomrule
\end{tabularx}
\end{table}

\begin{thebibliography}{9}
\bibitem{pinedo} M. Pinedo. \textit{Scheduling: Theory, Algorithms, and Systems}. Springer.
\bibitem{law} A. M. Law. \textit{Simulation Modeling and Analysis}. McGraw-Hill.
\bibitem{sa} S. Kirkpatrick, C. Gelatt, M. Vecchi. Optimization by Simulated Annealing. \textit{Science}, 1983.
\bibitem{tabu} F. Glover. Tabu Search. \textit{ORSA Journal on Computing}, 1989.
\bibitem{nsga2} K. Deb et al. A Fast and Elitist Multiobjective Genetic Algorithm: NSGA-II. \textit{IEEE Transactions on Evolutionary Computation}, 2002.
\bibitem{dqn} V. Mnih et al. Human-level Control through Deep Reinforcement Learning. \textit{Nature}, 2015.
\end{thebibliography}

\end{document}
